%%%%%%%%%%  scan idx 31  |  folio 18  %%%%%%%%%%
pos\'e est
logiquement ind\'ependant de tous les autres.

Enfin, dans Exp.\ XI, de nature nettement plus g\'en\'erale que XII et XIII et plus proche de l'esprit g\'en\'eral du S\'eminaire, on \'etudie, sur un topos localement annel\'e quelconque, un foncteur remarquable appel\'e «~\uline{foncteur d\'eterminant}~»:
\begin{equation}\tag{7.2}
\det^*:\underline{\operatorname{Parfis}}(X)\longrightarrow \underline{\Inv}(X)
\end{equation}
de la cat\'egorie des complexes parfaits sur $X$, les morphismes \'etant les seuls isomorphismes au sens de la cat\'egorie d\'eriv\'ee $D(X)$, dans la cat\'egorie des Modules inversibles sur $X$. Ce foncteur est d\'ecrit \`a peu de choses pr\`es par l'exigence d'\^etre de «~nature locale~», et d'\^etre donn\'e, pour un complexe $L$. born\'e et \`a composantes localement libres, par la formule
\begin{equation}\tag{7 3}
\det^*(L_\bullet)\simeq\bigotimes_i \det(L_i)^{(-1)^i},
\end{equation}
o\`u $\det(L_i)$ d\'esigne la puissance ext\'erieure de degr\'e maximum du Module localement libre $L_i$. Bien que cet expos\'e aussi soit logiquement ind\'ependant du reste du S\'eminaire \`a l'exception de l'expos\'e I, il constitue n\'eanmoins un \'el\'ement important du «~yoga~» g\'en\'eral propos\'e dans le S\'eminaire, et d\'evelopp\'e pour l'essentiel dans les expos\'es I \`a XI.

